\section{Verdikt}
\label{sec:verdikt}

Dieser Teil zieht die Schlussfolgerung aus den Befunden der
Teile~\ref{sec:unternehmen} bis \ref{sec:urteil}. Er f\"uhrt keine neue Zahl
ein; jede genannte Gr\"o{\ss}e steht bereits oben.

\vd{\textbf{Der rote Text ist eine Empfehlung, kein Beleg.}} Er ist aus den
vorstehenden Befunden unter den in Abschnitt~\ref{sec:annahmen} erkl\"arten
Annahmen abgeleitet. Schwarz bleibt, was belegt ist; blau sind
Einsch\"atzungen im laufenden Text; rot ist das Votum. Die Entscheidung trifft
der Verfasser.

\subsection{Votum}

\vd{\textbf{Nicht kaufen zum aktuellen Kurs von
\USDje{\FlowEinstiegskurs}. Beobachten und auf einen der Ausl\"oser in
Abschnitt~\ref{sec:ausloeser} warten.} Kein Verkaufsvotum: Wer die Aktie
h\"alt, hat aus dieser Analyse keinen zwingenden Grund zu verkaufen, denn die
Bilanz ist unbelastet, die Reserven wachsen und der Ausbau ist finanziert --
wohl aber einen Grund, die Position nicht aufzustocken. Wer sie nicht h\"alt,
hat keinen Grund einzusteigen, denn er bezahlt einen Preis, den keiner der
gerechneten F\"alle tr\"agt.}

\subsection{Begr\"undung aus den Befunden}
\label{sec:begruendung}

\vd{Vier Befunde tragen das Votum, und sie zeigen in verschiedene Richtungen.

\emph{Gegen einen Kauf spricht der Preis.} Der interne Zinsfu{\ss} erreicht in
keinem der drei Szenarien und \"uber keinen der drei Horizonte den
Ma{\ss}stab von \Pct{\FlowHuerde}; im oberen Szenario bleibt er \"uber
f\"unfzehn Jahre bei \Pct{\FlowIRRCFuenfzehn}
(Tabelle~\ref{tab:flowirr}). Das gilt auch dann noch, wenn man den
vorsichtigen Buchwert durch den vom Unternehmen selbst ver\"offentlichten
Kapitalwert des Island-Gold-Distrikts ersetzt
(Tabelle~\ref{tab:flowendwertvariante}).

\emph{Gegen einen Kauf spricht die Lage im Zyklus.} Der freie Cash-Flow des
Jahres 2025 ist der h\"ochste der Zehnjahresreihe und liegt im
\Pct{\FCFPerzentil}-Perzentil; der Median derselben Reihe betr\"agt
\MioUSDexakt{\FCFMedian}. Wer heute kauft, kauft auf dem Niveau des besten
Jahres, das dieses Unternehmen je hatte, bei den zweith\"ochsten
St\"uckkosten der Reihe.

\emph{F\"ur ein Beobachten statt eines Meidens spricht die Substanz.} Die
Reserven sind im \num{\ReservenJahreFolge}.\ Jahr in Folge gestiegen, zuletzt
um \Pct{\ReservenDelta} auf \MozReserven{\Reserven}; die Nettoliquidit\"at
betr\"agt \MioUSDexakt{\Nettoliquiditaet}; der ver\"offentlichte Ausbau bringt
\Pct{\ProgAchtProdDelta} mehr Produktion bei \Pct{\ProgAchtAISCRueckgang}
niedrigeren St\"uckkosten. Das ist kein Unternehmen, das man meidet, sondern
eines, dessen Preis man abwartet.

\emph{Und gegen das eigene Votum spricht der Konsens.} Die
\num{\KonsensZahl} sch\"atzenden H\"auser votieren mit \emph{Strong Buy}; das
mittlere Kursziel liegt \Pct{\KonsensAbstandKurs} \"uber dem heutigen Kurs,
und selbst das \emph{niedrigste} Ziel der gesamten Abdeckung liegt noch
\Pct{\KonsensTiefAbstandKurs} dar\"uber (Tabelle~\ref{tab:konsens}). Wenn eine
vollst\"andige Abdeckung oberhalb der eigenen Schwelle liegt, ist die
wahrscheinlichste Erkl\"arung eine eigene Annahme und kein gemeinsamer Irrtum
von \num{\KonsensZahl} H\"ausern. Der Mechanismus ist benennbar: Ein Kursziel auf
zw\"olf Monate bewertet den F\"orderer \"uber Gewinnvielfache und einen
unterstellten Fortf\"uhrungswert; eine Zahlungsreihe mit Ausstieg zum Buchwert
tut das nicht. Der Unterschied ist der Bewertungsaufschlag, nicht die
zugrunde liegende Zahl. Die Analysten rechnen zudem mit steigenden
Goldpreisen; dieses Dokument schreibt erzielte Preise konstant fort. Wer die
Erwartung steigender Preise teilt, kommt zu ihrem und nicht zu diesem
Ergebnis.}

\subsection{Einstiegsschwelle}
\label{sec:schwelle}

Tabelle~\ref{tab:schwelle} nennt den Kurs, bei dem der interne Zinsfu{\ss} den
Ma{\ss}stab von \Pct{\FlowHuerde} genau trifft. Das ist die Zahl, die eine
Anlageentscheidung braucht: nicht ob, sondern bis zu welchem Preis.

\begin{table}[htbp]
\centering
\caption{Einstiegskurs in USD, bei dem der interne Zinsfu{\ss} die H\"urde von
\Pct{\FlowHuerde} erreicht. Endwert: Buchwert je Aktie
\USDje{\FlowEndwert}. Aktueller Kurs: \USDje{\FlowEinstiegskurs}.}
\label{tab:schwelle}
\begin{tabular}{lrrrr}
\toprule
Szenario & Goldpreis & 5 Jahre & 10 Jahre & 15 Jahre\\
\midrule
\tabellenkoerper{data/flow_schwelle}
\bottomrule
\end{tabular}
\end{table}

\FloatBarrier

\vd{Die Schwelle liegt im mittleren Szenario \"uber f\"unfzehn Jahre bei
\USDje{\FlowSchwelleBasis}; der heutige Kurs liegt \Pct{\FlowAbstandBasis}
dar\"uber. Im oberen Szenario -- dem h\"ochsten je erzielten Preis, dauerhaft
fortgeschrieben -- steigt sie auf \USDje{\FlowSchwelleBeste}, und der Abstand
schrumpft auf \Pct{\FlowAbstandBeste}. Mit dem g\"unstigsten belegbaren
Endwert liegt sie bei \USDje{\FlowSchwelleIGDBeste}, und der Abstand
betr\"agt immer noch \Pct{\FlowAbstandIGDBeste}.

Das ist der Kern des Falls: Es gibt keine Kombination aus einem vom
Unternehmen berichteten Goldpreis, einem belegten Endwert und einem der drei
Horizonte, bei der der heutige Kurs die eigene H\"urde des Unternehmens
tr\"agt.}

\subsubsection{Die Wachstumsrate, die der Preis voraussetzt}

Den Cash-Flow konstant zu halten h\"alt die Rechnung ehrlich und l\"asst
zugleich jedes wachsende Unternehmen teuer aussehen. Statt das durch eine
Prognose zu heilen, kehrt Tabelle~\ref{tab:wachstum} den Preis in die
Wachstumsrate um, die er bereits voraussetzt.

\begin{table}[htbp]
\centering
\caption{J\"ahrliche Wachstumsrate des verteilbaren Cash-Flows in Prozent, die
der heutige Kurs von \USDje{\FlowEinstiegskurs} voraussetzt, damit der interne
Zinsfu{\ss} \Pct{\FlowHuerde} erreicht.}
\label{tab:wachstum}
\begin{tabular}{lrrr}
\toprule
Szenario & 5 Jahre & 10 Jahre & 15 Jahre\\
\midrule
\tabellenkoerper{data/flow_wachstum}
\bottomrule
\end{tabular}
\end{table}

\FloatBarrier

\vd{Damit wird aus einem Urteil eine pr\"ufbare Behauptung. Der heutige Kurs
enth\"alt im mittleren Szenario \Pct{\FlowWachstumBZehn} Wachstum des
verteilbaren Cash-Flows \emph{jedes Jahr, zehn Jahre lang} -- zus\"atzlich zu
dem Ausbau, den die Rechnung bereits enth\"alt. \"Uber f\"unfzehn Jahre sind es
\Pct{\FlowWachstumBFuenfzehn}. Dieses Unternehmen hat \"uber die letzten zehn
Jahre einen Median des freien Cash-Flows von \MioUSDexakt{\FCFMedian}
geliefert (Tabelle~\ref{tab:reihe}). Eine solche Rate ist nicht unm\"oglich --
sie ist bei einem Goldproduzenten vor allem eine Aussage \"uber den Goldpreis
-- aber sie steht im Preis, und der K\"aufer bezahlt sie im Voraus.}

\subsection{Ausl\"oser, die das Votum drehen}
\label{sec:ausloeser}

\vd{Alle sind kostenlos zu beobachten; zwei werden quartalsweise
ver\"offentlicht, einer steht jeden Tag im Kurszettel.

\emph{Zum Kauf.} Der Kurs f\"allt unter \USDje{\FlowSchwelleBeste}, ohne dass
sich fundamental etwas verschlechtert; dann tr\"agt das obere Szenario. Oder
die Produktion erreicht \"uber zwei aufeinanderfolgende Quartale den
ver\"offentlichten Prognosepfad von \TsdUnzen{\FlowMengeSechs} auf
Jahresbasis, \emph{und} die AISC fallen wieder in die urspr\"ungliche Spanne
von \JeUnze{\PrognoseAiscUnten} bis \JeUnze{\PrognoseAiscOben} zur\"uck. Dann
ist die Kostenanhebung vom Juli 2026 eine Episode gewesen und kein Trend, und
der Abstand zum Sektor aus Abschnitt~\ref{sec:zerlegung} sollte sich
schlie{\ss}en.

\emph{Zum Meiden.} Die Fertigstellung der Schachterweiterung verschiebt sich
\"uber das Jahresende 2026 hinaus, oder der Dreijahresausblick f\"ur 2028 wird
gesenkt. Beides tr\"afe die Gr\"o{\ss}e, auf der die ganze Rechnung ab 2028
ruht. Ebenfalls zum Meiden: ein R\"uckgang des erzielten Goldpreises in die
N\"ahe des Standes von 2024, denn dann tr\"agt der laufende Cash-Flow das
Bauprogramm nicht mehr (Abschnitt~\ref{sec:kapazitaet}) und die
unbeanspruchte Kreditlinie wird zur Finanzierungsquelle.}

\subsection{Was dieses Votum nicht wissen kann}

\vd{Die entscheidende Gr\"o{\ss}e ist nicht berechenbar: der k\"unftige
Goldpreis. Tabelle~\ref{tab:schwelle} sagt, was jede Annahme wert ist; welche
zutrifft, sagt sie nicht. Dieses Dokument schreibt erzielte Preise konstant
fort, weil es das Einzige ist, was sich belegen l\"asst -- und genau das ist
bei einem Rohstoffproduzenten die schw\"achste Stelle der Methode.

Drei weitere Vorbehalte, die gegen das eigene Votum sprechen.

\emph{Der Konsens steht dagegen, und zwar geschlossen.} Kein einziges
ver\"offentlichtes Kursziel liegt unter dem heutigen Kurs
(Abschnitt~\ref{sec:begruendung}). Eine vollst\"andige Abdeckung auf der
Gegenseite ist ein Grund, die eigene Endwertannahme f\"ur zu vorsichtig zu
halten, und nicht bloss eine Fu{\ss}note.

\emph{Der Ausbau ist gr\"o{\ss}er als das, was die Rechnung erfasst.} Die
Rechnung nimmt die Produktion des Dreijahresausblicks bis 2028 auf und
h\"alt sie danach konstant. Sie enth\"alt damit weder Lynn Lake \"uber das
Jahr 2029 hinaus in voller Gr\"o{\ss}e noch PDA noch weitere Umwandlungen von
Ressourcen in Reserven -- und die Reserven sind im
\num{\ReservenJahreFolge}.\ Jahr in Folge gestiegen. Trifft dieses Muster
weiter zu, ist der ausgewiesene interne Zinsfu{\ss} eher eine Untergrenze.

\emph{Und der Zeitpunkt k\"onnte bereits g\"unstig sein.} Die Aktie steht
\PctBetrag{\KursAbstandHoch} unter ihrem Hoch vom \EinbruchTagHoch{} und hat
die Erholung ihres Sektors um \Pct{\AbstandUeberschuss} nicht mitgemacht. Wer
das Votum ablehnt, kann mit gutem Grund einwenden, dass genau dieser Abstand
die Gelegenheit ist, auf die es zu warten vorgibt.}
